% Compile with pdfLaTeX. Revised after Majid Zamani's October 23 remaining comments.
% All degree bounds mean total degree at most the prescribed d.
\documentclass[10pt,letterpaper]{article}
\usepackage[T1]{fontenc}
\usepackage{lmodern}
\usepackage[margin=0.8in]{geometry}
\usepackage{amsmath,amssymb,amsthm,booktabs,array,microtype}
\usepackage{enumitem,needspace}
\usepackage[dvipsnames]{xcolor}
\usepackage[colorlinks=true,linkcolor=MidnightBlue,urlcolor=MidnightBlue,citecolor=MidnightBlue]{hyperref}
\newtheorem{theorem}{Theorem}
\newtheorem{lemma}{Lemma}
\newtheorem{proposition}{Proposition}
\newtheorem{corollary}{Corollary}
\theoremstyle{remark}
\newtheorem*{remark}{Remark}
\newcommand{\R}{\mathbb R}
\newcommand{\co}{\operatorname{conv}}
\newcommand{\cl}{\operatorname{cl}}
\newcommand{\dist}{\operatorname{dist}}
\newcommand{\one}{\mathbf 1}
\newcommand{\ball}{\mathbb B_V}
\newcommand{\F}{\mathcal F_d}
\newcommand{\K}{\mathcal K_\varepsilon}
\newcommand{\dd}{\,\mathrm d}
\hypersetup{pdftitle={Schur-Stable Fixed-Anchor Certificates at Prescribed Degree},
 pdfauthor={Reza Iraji},
 pdfsubject={Robust feature-space converse and certified finite construction}}
\setlength{\parindent}{0pt}
\setlength{\parskip}{3pt}
\setlist[enumerate]{leftmargin=*,itemsep=2pt,topsep=3pt}
\setlist[itemize]{leftmargin=*,itemsep=2pt,topsep=3pt}
\title{Schur-Stable Fixed-Anchor Certificates at Prescribed Degree:\\A Robust Converse and Conditional Finite Extraction}
\author{Reza Iraji\\{\small Technical Note 3, revised after the October 23 comments}}
\date{October 25, 2026}
\begin{document}
\maketitle
\vspace{-1em}
\begin{abstract}
At a prescribed polynomial degree, robust separation of the convexified monomial dynamics is equivalent to a finite forward fixed-anchor vector barrier certificate with a constant nonnegative Schur-stable comparison matrix. This is exactly the class that can be shifted to have a uniform strict global propagation margin. It forces uniform eventual avoidance of the unsafe set from the entire verification domain, and therefore excludes onto maps and unsafe periodic orbits. Conditional extraction from a supplied, certified convex body gives an explicit component bound, together with a finite linear program that accounts for source and successor approximation errors. A nonlinear worked certificate is replayed with exact rational arithmetic and compared with its simpler scalar baseline. A separate perturbation-stable nonlinear family has minimum scalar comparison degree two and minimum finite-vector degree one. The moment obstruction and exact implication-versus-comparison degree separation are retained with their precise scopes; physical robust safety alone does not imply the feature-space hypothesis.
\end{abstract}

\section{Setting and the convexified dynamics}
Fix an integer $d\geq1$. Assume $X\subseteq\R^n$ is nonempty and compact, $f:X\to X$ is continuous, and $X_0,X_u\subseteq X$ are nonempty, compact, and disjoint. All degree bounds below mean \emph{total degree at most $d$}, with $d$ chosen in advance. A forward fixed-anchor VBC consists of finitely many polynomial components $B=(B_1,\ldots,B_m)^\top$, a constant entrywise nonnegative $A$, and $\delta_0,\delta_u>0$ such that
\begin{equation}\label{eq:ordinary}
 B\leq-\delta_0\one\ \text{on }X_0,\qquad
 B_1\geq\delta_u\ \text{on }X_u,\qquad
 B(f(x))\leq AB(x)\quad(x\in X).
\end{equation}
The same designated component $B_1$ separates \emph{all} unsafe states. No restriction on $\rho(A)$ is assumed. The converse in this note concerns the forward direction.

Let $\phi_d(x)\in\R^D$, $D=\binom{n+d}{d}-1$, list the nonconstant monomials of degree at most $d$. Constants are represented by affine offsets, so every admissible polynomial is $B_i(x)=\beta_i(\phi_d(x))$ for an affine $\beta_i$. Define
\begin{align}
 E&=\operatorname{aff}\phi_d(X),\quad V=E-E,\quad
 Z=\co\phi_d(X),\nonumber\\
 I_d&=\co\phi_d(X_0),\qquad J_d=\co\phi_d(X_u),\label{eq:features}\\
 \F&=\co\{(\phi_d(x),\phi_d(f(x))):x\in X\}\subseteq E\times E.\nonumber
\end{align}
The convex sets $Z$, $I_d$, $J_d$, and $\F$ are compact; the affine space $E$ and direction space $V$ need not be bounded. The unit ball $\ball$ is in $V$, and interiors, distances, and norms are relative to $E$. Omitting the constant feature avoids artificial perturbations of its identically-one coordinate.

A pair $(s,t)\in\F$ can be the feature averages of a distribution and its successor distribution, even when $s$ and $t$ are not individual state lifts. An affine comparison on the original graph holds on all of $\F$. Thus ordinary lifted trajectory safety does not suffice.

For $\varepsilon\geq0$, put
\begin{align}
 \mathcal F_d^\varepsilon&=\F+(\varepsilon\ball\times\varepsilon\ball),\nonumber\\
 T_\varepsilon(C)&=\{t:\ \exists s\in C,\ (s,t)\in\mathcal F_d^\varepsilon\},\nonumber\\
 C_\varepsilon^0&=I_d+\varepsilon\ball,\qquad
 C_\varepsilon^{k+1}=\co\bigl(C_\varepsilon^k\cup T_\varepsilon(C_\varepsilon^k)\bigr),\label{eq:closure}\\
 \mathcal K_\varepsilon&=\cl\bigcup_{k\geq0}C_\varepsilon^k.\nonumber
\end{align}
The source and successor errors are independent feature-space errors, not a claim about physical state disturbances. There are two convexifications: the transition graph in $\F$, and the reachable set at every iteration in \eqref{eq:closure}. Each $C_\varepsilon^k$ is compact and convex; the sets increase and lie in the compact set $Z+\varepsilon\ball$. Consequently $\mathcal K_\varepsilon$ is compact and convex.

Every VBC satisfying \eqref{eq:ordinary} forces
\begin{equation}\label{eq:necessary}
 \mathcal K_0\cap J_d=\varnothing.
\end{equation}
Indeed, $\{\beta\leq0\}$ contains $I_d$ and is preserved by $\F$, convex combinations, and limits, whereas $\beta_1\geq\delta_u$ on $J_d$. This is a necessary condition; no general converse to \eqref{eq:necessary} is asserted here.

\section{The converse and its exact comparison-matrix class}
\begin{theorem}[Robust convexified safety, strict propagation, and Schur stability]\label{th:main}
Under the assumptions above, the following are equivalent.
\begin{enumerate}
\item There is an $\varepsilon>0$ such that $\mathcal K_\varepsilon\cap J_d=\varnothing$.
\item There are finite $m\geq1$, polynomial $B_i$ of degree at most $d$, a constant $A\geq0$, and $\delta_0,\delta_u,\eta>0$ satisfying the separation conditions in \eqref{eq:ordinary} and
\begin{equation}\label{eq:strict}
 B(f(x))\leq AB(x)-\eta\one\qquad(x\in X).
\end{equation}
\item There is a VBC satisfying \eqref{eq:ordinary}, at the prescribed degree bound, with a constant $A\geq0$ satisfying $\rho(A)<1$.
\end{enumerate}
Every matrix in item 2 is necessarily Schur stable. Items 2 and 3 can be interconverted without changing the degree or component count, although item 1 does not specify a minimal count.
Moreover, the robust separation condition in item 1 yields a bounded polyhedron
\[
 P=\{z\in E:\beta_j(z)\leq0,\ j=1,\ldots,m\}
\]
with $I_d$ strictly inside it, one designated supporting inequality strictly separating $J_d$, and
\[
 (s,t)\in\F,\ s\in P\quad\Longrightarrow\quad t+\varepsilon\ball\subseteq P.
\]
Thus it yields a VBC of the prescribed degree bound in the usual class \eqref{eq:ordinary}.
\end{theorem}

\textbf{Scope.} The equivalence characterizes the \emph{Schur-stable subclass} of ordinary forward fixed-anchor VBCs, rather than all certificates in \eqref{eq:ordinary}. Ordinary certificates can have $\rho(A)=1$. The converse does not cover that boundary class. The existence proof below uses compactness and continuity; Sections~\ref{sec:quant}--\ref{sec:lp} add finite information and explicit approximation guarantees.

\subsection{Two finite-dimensional ingredients}
\begin{lemma}[Finite outer approximation with a prescribed separator]\label{lem:poly}
Let $K\subseteq E$ be a compact convex body, let $a_1\in V$ have unit norm, and let $\tau>0$. There are finitely many unit vectors $a_1,\ldots,a_m\in V$ such that, with $h_K(a)=\max_{z\in K}\langle a,z\rangle$,
\[
 K\subseteq P=\{z\in E:\langle a_j,z\rangle\leq h_K(a_j),\ j=1,\ldots,m\}
 \subseteq K+\tau\ball,
\]
and $P$ is bounded.
\end{lemma}
\begin{proof}
Choose an orthonormal basis of $V$ and its positive and negative coordinate directions. Their supporting halfspaces bound a compact box $Q$ containing $K$. Each $y\in Q$ with $\dist(y,K)\geq\tau$ is strictly excluded by a supporting halfspace: if $p$ is its closest point in $K$, take $a=(y-p)/\|y-p\|$, so $\langle a,y\rangle>h_K(a)=\langle a,p\rangle$. The corresponding open exclusion sets cover the compact set $\{y\in Q:\dist(y,K)\geq\tau\}$. A finite subcover, the box inequalities, and the prescribed direction $a_1$ give the result.
\end{proof}

\Needspace{9\baselineskip}
\begin{lemma}[Global row multiplier from a convex relation]\label{lem:dual}
Let $F\subseteq E\times E$ be compact and convex, let $\beta_1,\ldots,\beta_m$ be affine, and suppose some $(s_*,t_*)\in F$ has $\beta(s_*)<0$. If
\[
 p_i=\max\{\beta_i(t):(s,t)\in F,\ \beta(s)\leq0\},
\]
then there is an attained $a_i\in\R_{\geq0}^m$ such that
\[
 \beta_i(t)-a_i^\top\beta(s)\leq p_i\quad\text{for every }(s,t)\in F.
\]
\end{lemma}
\begin{proof}
The closed convex set
\[
 H_i=\{(\beta(s)+u,\ \beta_i(t)-v):(s,t)\in F,\ u\geq0,\ v\geq0\}
\]
contains $(0,p_i)$ on its boundary. It is closed because it is a compact affine image plus a closed cone. A nonzero supporting normal $(w,\gamma)$ can be oriented so that
\[
 w^\top y+\gamma q\leq\gamma p_i\quad((y,q)\in H_i).
\]
Then $w\leq0$ and $\gamma\geq0$, since $u$ and $v$ can increase without bound. If $\gamma=0$, strict feasibility gives $w^\top\beta(s_*)>0$ for nonzero $w\leq0$, contradicting the displayed inequality. Hence $\gamma>0$; dividing by $\gamma$ and taking $a_i=-w/\gamma$ proves the claim. This is the usual affine-constraint convex duality argument.
\end{proof}

\subsection{Proof of Theorem~\ref{th:main}}
\begin{proof}
\emph{(1)$\Rightarrow$(2): from the closure to finitely many facets.}
Write $K=\mathcal K_\varepsilon$. It has nonempty relative interior because it contains $I_d+\varepsilon\ball$. Compact convex separation of $K$ and $J_d$ provides a unit $a_1\in V$ and
\begin{equation}\label{eq:anchor}
 \gamma=\min_{z\in J_d}\langle a_1,z\rangle-h_K(a_1)>0.
\end{equation}
Apply Lemma~\ref{lem:poly} with $\tau=\varepsilon/2$, and define
\[
 \beta_j(z)=\langle a_j,z\rangle-h_K(a_j),\qquad
 B_j(x)=\beta_j(\phi_d(x)).
\]
Each polynomial has degree at most $d$. Since $I_d+\varepsilon\ball\subseteq K$ and every $a_j$ is a unit vector,
\[
 \beta_j(s)\leq-\varepsilon\quad(s\in I_d),\qquad
 \beta_1(z)\geq\gamma\quad(z\in J_d).
\]

The key transition fact is
\begin{equation}\label{eq:key}
 (s,t)\in\F,\ \dist(s,K)\leq\varepsilon/2
 \quad\Longrightarrow\quad t+\varepsilon\ball\subseteq K.
\end{equation}
To prove it, choose $p\in K$ closest to $s$, and $p_\ell\in\bigcup_k C_\varepsilon^k$ tending to $p$. For large $\ell$, $\|p_\ell-s\|<\varepsilon$. For every $w\in\varepsilon\ball$, the pair $(p_\ell,t+w)$ belongs to $\mathcal F_d^\varepsilon$. If $p_\ell\in C_\varepsilon^{k_\ell}$, then $t+w\in C_\varepsilon^{k_\ell+1}\subseteq K$. This proves \eqref{eq:key}. In particular, the proof does not assume that passing to the closure automatically preserves every boundary source fiber of the perturbed relation.

For $(s,t)\in\F$ with $s\in P$, the outer approximation gives $\dist(s,K)\leq\varepsilon/2$. By \eqref{eq:key}, $t+\varepsilon\ball\subseteq K$, hence
\[
 \beta_i(t)\leq-\varepsilon\quad(i=1,\ldots,m).
\]
There is a strictly feasible source in $\F$, obtained from any $x_0\in X_0$. Lemma~\ref{lem:dual}, applied row by row, therefore supplies $A_i\geq0$ with
\[
 \beta_i(t)-A_i\beta(s)\leq-\varepsilon\quad((s,t)\in\F).
\]
Stacking the rows and restricting to the original graph yields \eqref{eq:strict} with $\eta=\delta_0=\varepsilon$ and $\delta_u=\gamma$. The same argument gives the stated relational invariance of $P$.

\emph{(2)$\Rightarrow$(1): strict propagation tolerates feature errors.}
Write the linear part of $\beta_i$ on $V$ as $a_i$, let $L_i=\|a_i\|$, and set
\[
 L=\max_i L_i>0,\qquad
 M=\max_i\left(L_i+\sum_j A_{ij}L_j\right)>0.
\]
The positivity follows from the separating anchor. Choose
\[
 0<\varepsilon<\min\{\delta_0/(2L),\ \eta/(2M)\}.
\]
Affine linearity extends \eqref{eq:strict} from the graph to $\F$. For $(s,t)\in\F$ and $e,w\in\varepsilon\ball$,
\begin{align*}
 \beta_i(t+w)
 &\leq (A\beta(s+e))_i-\eta
       +\varepsilon\left(L_i+\sum_j A_{ij}L_j\right)\\
 &\leq (A\beta(s+e))_i-\eta/2.
\end{align*}
Also $\beta\leq-\delta_0\one/2$ on $C_\varepsilon^0$. Nonnegativity of $A$ and the displayed perturbed comparison imply inductively $\beta\leq0$ on every $C_\varepsilon^k$; affine inequalities persist under convex combinations and limits. Thus $\beta\leq0$ on $\mathcal K_\varepsilon$, whereas $\beta_1\geq\delta_u$ on $J_d$. The sets are disjoint.

\emph{(2)$\Rightarrow$(3): strict propagation forces Schur stability.}
For a nonzero left Perron vector $w\geq0$, write $w^\top A=\rho(A)w^\top$. Fix $x_0\in X_0$ and $q_k=w^\top B(f^k(x_0))$. Then $q_0<0$ and
\[
 q_{k+1}\leq\rho(A)q_k-\eta w^\top\one.
\]
If $\rho(A)\geq1$, induction makes this upper bound tend to $-\infty$, contradicting boundedness of $B$ on the compact forward-invariant $X$. Thus $\rho(A)<1$. The same argument without $\eta$ excludes $\rho(A)>1$ for every ordinary VBC in \eqref{eq:ordinary}; consequently such a VBC always has $\rho(A)\leq1$.

\emph{(3)$\Rightarrow$(2): a constant shift supplies strictness.}
Let $v=(I-A)^{-1}\one=\sum_{j\geq0}A^j\one\geq\one$. Choose $0<c<\delta_u/v_1$ and put $\widetilde B=B-cv$. Since $v-Av=\one$,
\[
 \widetilde B\circ f\leq A\widetilde B-c\one,\qquad
 \widetilde B\leq-\delta_0\one-cv\ \text{on }X_0,\qquad
 \widetilde B_1\geq\delta_u-cv_1>0\ \text{on }X_u.
\]
This preserves the degree and number of components.
\end{proof}

\begin{remark}[Rational approximation]
Strict margins on compact $X$ also imply existence of a VBC with rational polynomial coefficients and a rational matrix $A\geq0$, at the same degree and component count. Indeed, the separation values and the residual $AB-B\circ f$ depend continuously on these finitely many coefficients, uniformly on $X$, since $f(X)\subseteq X$. Sufficiently close rational approximations, keeping $A\geq0$ and errors below half the respective positive margins, preserve all inequalities. This is an existence guarantee, not a computable rounding procedure for unspecified compact data.
\end{remark}

\begin{corollary}[A sufficient condition that can be checked on a chosen body]\label{cor:body}
It suffices to exhibit a compact convex $C\subseteq E$ and $\varepsilon>0$ satisfying
\[
 I_d+\varepsilon\ball\subseteq C,\qquad
 T_\varepsilon(C)\subseteq C,\qquad C\cap J_d=\varnothing.
\]
Then $\mathcal K_\varepsilon\subseteq C$, and Theorem~\ref{th:main} constructs a finite degree-$d$ fixed-anchor VBC.
\end{corollary}
This is a robust convex-body assumption, not an assumption that a finite polyhedral invariant already exists. The theorem uses supporting halfspaces to obtain finitely many functions and the global comparison matrix.

\begin{remark}[Finite continuous features]
The proof uses only compactness of the graph image and affinity of the component functions in feature space. It therefore applies to any fixed finite continuous feature map, with affine offsets representing constants and with all balls relative to its affine hull. The monomial choice supplies the prescribed-polynomial-degree interpretation. No computational claim follows from continuity alone.
\end{remark}

\section{Dynamical consequences and scope limitations}\label{sec:limitation}
\begin{proposition}[Uniform eventual avoidance on the entire domain]\label{prop:absorb}
Every strict certificate in Theorem~\ref{th:main} has a common integer $N$ such that $B(f^k(x))<0$ componentwise for every $x\in X$ and $k\geq N$. Thus all trajectories in the entire verification domain eventually avoid $X_u$, uniformly in their starting point. The same avoidance conclusion holds for ordinary Schur-stable certificates, using their shifted functions.
\end{proposition}
\begin{proof}
Iteration gives
\[
 B(f^k(x))\leq A^kB(x)-\eta\sum_{j=0}^{k-1}A^j\one
 =A^k(B(x)+\eta v)-\eta v,\quad v=(I-A)^{-1}\one.
\]
Since $\rho(A)<1$ and $B$ is bounded, the right side tends uniformly to $-\eta v<0$. More quantitatively, let
\[
 \kappa=\max_i\frac{(Av)_i}{v_i}=\max_i(1-v_i^{-1})<1,\qquad
 H=\max\left\{0,\ \max_{x\in X,i}\frac{B_i(x)+\eta v_i}{v_i}\right\}.
\]
Nonnegativity of $A$ yields $B(f^k(x))\leq(H\kappa^k-\eta)v$. Any $N$ with $H\kappa^N\leq\eta/2$ suffices. If $\kappa=0$, take $N=1$; if $H=0$, take $N=0$. For an ordinary Schur-stable certificate apply this argument to $\widetilde B$ from the theorem; the original unshifted functions need not become negative.
\end{proof}

\begin{corollary}[Obstructions at every degree]\label{cor:onto}
An unsafe periodic point excludes strict and ordinary Schur-stable fixed-anchor VBCs at every degree and finite component count. If $f$ is onto $X$, these classes are also excluded whenever $X_u\neq\varnothing$.
\end{corollary}
\begin{proof}
An unsafe periodic point is visited at arbitrarily late times, contradicting Proposition~\ref{prop:absorb}. If $f$ is onto, $f^N(X)=X$; uniform eventual negativity would force the anchor to be negative on all of $X$, including $X_u$.
\end{proof}
These are limitations of the additional strict/Schur-stable class, not obstructions to all ordinary certificates.

\section{A quantitative finite-facet construction}\label{sec:quant}
We now strengthen the chosen-body criterion (Corollary~\ref{cor:body}). The body is an input to be certified, not an unknown infinite-horizon closure whose support values are assumed available. Fix $z_c\in E$, let $r=\dim V\geq1$, and write
\[
 h_C(a)=\max_{z\in C}\langle a,z-z_c\rangle.
\]
Assume a compact convex body $C$ and certified constants $\varepsilon,\gamma,R>0$, with $R\geq\varepsilon$, satisfy
\begin{equation}\label{eq:bodydata}
 I_d+\varepsilon\ball\subseteq C,\quad T_\varepsilon(C)\subseteq C,\quad
 C\subseteq z_c+R\ball,\quad
 \min_{z\in J_d}\langle a_1,z-z_c\rangle\geq h_C(a_1)+\gamma,
\end{equation}
where $a_1\in V$ is a certified unit anchor direction.

\begin{lemma}[A coarse explicit support net]\label{lem:net}
Set $\delta=\varepsilon/(16R)$ and choose
$0\leq\alpha\leq\min\{\varepsilon/8,\gamma/2\}$.
There is a finite $\delta$-net $Q$ of the unit sphere in $V$, in Euclidean chord distance, with
\[
 |Q|\leq\left\lceil(1+2/\delta)^r-(2/\delta-1)^r\right\rceil.
\]
Use the directions $Q\cup\{a_1\}$, with the anchor first, and certified support upper bounds
\[
 h_C(a_i)\leq\widehat h_i\leq h_C(a_i)+\alpha.
\]
For $\beta_i(z)=\langle a_i,z-z_c\rangle-\widehat h_i$ and $P=\{\beta\leq0\}$,
\[
 C\subseteq P\subseteq C+\tfrac{\varepsilon}{2}\ball.
\]
\end{lemma}
\begin{proof}
Choose a maximal $\delta$-separated subset of the sphere. The disjoint open balls of radius $\delta/2$ about its points lie in the annulus with radii $1-\delta/2$ and $1+\delta/2$. Comparing $r$-dimensional volumes proves the cardinality bound and ensures the selection is finite. Maximality gives the net property.

For $z\in P$, let $u=(z-z_c)/\|z-z_c\|$ when $z\neq z_c$ and take $q\in Q$ with $\|q-u\|\leq\delta$. As $\langle q,u\rangle\geq1-\delta$, we have
\[
 \|z-z_c\|\leq (R+\alpha)/(1-\delta).
\]
For every unit $a$ take a net direction $q$ within $\delta$. Since the support function is $R$-Lipschitz on the unit sphere,
\begin{align*}
 \langle a,z-z_c\rangle-h_C(a)
 &\leq\alpha+\delta\|z-z_c\|+R\delta\\
 &\leq 2\alpha+4R\delta\leq\varepsilon/2.
\end{align*}
Here $\delta\leq1/16$, so $(1-\delta)^{-1}\leq16/15$ suffices for the second bound. The support-halfspace description of $C+(\varepsilon/2)\ball$ proves the inclusion. Lower support bounds give $C\subseteq P$.
\end{proof}

\begin{theorem}[Explicit component bound from a certified body]\label{th:count}
Under \eqref{eq:bodydata} and the support-value assumptions of Lemma~\ref{lem:net}, a strict forward fixed-anchor VBC exists with
\begin{equation}\label{eq:mcount}
 m\leq 1+\left\lceil
 (1+32R/\varepsilon)^r-(32R/\varepsilon-1)^r
 \right\rceil,
\end{equation}
and margins $\delta_0\geq\varepsilon$, $\delta_u\geq\gamma-\alpha\geq\gamma/2$, and $\eta\geq\varepsilon$. Its matrix is necessarily Schur stable. The geometric count is $O_r((R/\varepsilon)^{r-1})$; no optimality is claimed.
\end{theorem}
\begin{proof}
Use Lemma~\ref{lem:net}. The initial tube gives $\beta_i\leq-\varepsilon$ on $I_d$, and the anchor gives the stated unsafe margin. For a nominal pair $(s,t)\in\F$ with $\dist(s,C)\leq\varepsilon$, choose a closest $y\in C$. Then
\[
 (y,t+w)\in\mathcal F_d^\varepsilon\quad(w\in\varepsilon\ball),
 \qquad t+\varepsilon\ball\subseteq C.
\]
This uses the certified inclusion $T_\varepsilon(C)\subseteq C$ directly. In particular it holds for $s\in P$, so $\beta_i(t)\leq-\varepsilon$. Lemma~\ref{lem:dual} supplies the global nonnegative comparison rows. Theorem~\ref{th:main} gives Schur stability. The bound counts the net and at most one additional anchor direction.
\end{proof}
The unsafe margin $\gamma$ controls the required support accuracy through $\alpha$. It does not enlarge the count because its separating direction is supplied explicitly.

The same row argument works for \emph{any} certified unit-normal outer polyhedron satisfying
\begin{equation}\label{eq:otherpoly}
 C\subseteq P\subseteq C+\varepsilon\ball,
\end{equation}
with the stated support errors and anchor. Thus a more efficient polyhedron can replace the coarse universal net. Supporting-halfspace approximation and preservation of invariance within a dynamical margin are classical mechanisms; the count here is a quantitative consequence for this degree-restricted certificate class.

\subsection{A sharper count using certified projections}
The direct support net above gives a simple bound, but it is not the sharpest geometric route. The classical Dudley projection construction achieves the square-root dependence on approximation accuracy; see the self-contained proof by Har-Peled--Jones~\cite{harpeled}. The following variant tracks support errors explicitly before extracting the VBC.

\begin{proposition}[Projection-based component bound]\label{prop:projection}
Assume \eqref{eq:bodydata}, and assume exact Euclidean projections onto $C$ can be computed. Set
\[
 \ell=\tfrac12\sqrt{R\varepsilon},\qquad
 0\leq\alpha\leq\min\{\varepsilon/16,\gamma/2\}.
\]
Take an $\ell$-net of the sphere $\{q:\|q-z_c\|=2R\}$, obtained from a maximal $\ell$-separated set. At each net point $q_i$ compute $p_i=\operatorname{proj}_C(q_i)$ and the unit normal
$\nu_i=(q_i-p_i)/\|q_i-p_i\|$. Include these normals together with the given unsafe anchor, numbered first, and use support upper offsets within $\alpha$ of the true values. This constructs a polyhedron
\[
 C\subseteq P\subseteq C+(\varepsilon/2)\ball
\]
and hence a strict degree-at-most-$d$ fixed-anchor VBC with
\begin{equation}\label{eq:projection-count}
 m\leq1+\left\lceil
 (1+8\sqrt{R/\varepsilon})^r-
 (8\sqrt{R/\varepsilon}-1)^r
 \right\rceil
 =O_r\bigl((R/\varepsilon)^{(r-1)/2}\bigr).
\end{equation}
The initial, unsafe-anchor, and nominal propagation margins are respectively at least $\varepsilon$, $\gamma-\alpha$, and $\varepsilon$.
\end{proposition}
\begin{proof}
Translate $z_c$ to the origin. The projection property gives
$h_C(\nu_i)=\langle \nu_i,p_i\rangle$; thus every halfspace contains $C$.
Packing disjoint balls of radius $\ell/2$ in the spherical annulus of radii $2R-\ell/2$ and $2R+\ell/2$ gives at most
\[
 \left\lceil(1+4R/\ell)^r-(4R/\ell-1)^r\right\rceil
\]
net points. Adding the anchor proves the stated count.

We prove the outer-distance bound for every $z\in P\setminus C$.
Write $p=\operatorname{proj}_C(z)$, $t=\|z-p\|>0$, and $n=(z-p)/t$.
The ray $p+\lambda n$ meets the sphere of radius $2R$ at
$q=p+Ln$ with $R\leq L\leq3R$.
Choose a net point $q_i=q+e$ with $\|e\|\leq\ell$, put $d_i=p_i-p$, and
$u_i=q_i-p_i=Ln+e-d_i$.
The supporting-normal property at $p$ gives $\langle n,d_i\rangle\leq0$.
Consequently,
\begin{align*}
 \langle u_i,z-p_i\rangle
 &=Lt+t\langle n,e\rangle-(L+t)\langle n,d_i\rangle
       -\langle e,d_i\rangle+\|d_i\|^2\\
 &\geq (R-\ell)t-\ell^2/4.
\end{align*}
The last step uses
$\|d_i\|^2-\langle e,d_i\rangle
=\|d_i-e/2\|^2-\|e\|^2/4\geq-\ell^2/4$.
The facet with support overestimate at most $\alpha$ implies
$\langle u_i,z-p_i\rangle\leq\alpha\|u_i\|\leq3R\alpha$,
since $\|q_i\|=2R$ and $\|p_i\|\leq R$.
As $\ell\leq R/2$,
\[
 t\leq\frac{\ell^2/4+3R\alpha}{R-\ell}
 \leq\frac{\ell^2}{2R}+6\alpha
 \leq\frac{\varepsilon}{8}+\frac{6\varepsilon}{16}
 =\frac{\varepsilon}{2}.
\]
This proves boundedness and the outer inclusion. The initial/unsafe margin calculation and global row-duality argument from Theorem~\ref{th:count} now give the VBC and its stated margins.
\end{proof}
The improved exponent is the classical Dudley approximation rate, not an independent novelty claim.

\textbf{Certified inexact projection data.}
Exact projections are sufficient but can be relaxed. It is enough to provide feasible points $p_i\in C$ with
\[
 \nu_i=\frac{q_i-p_i}{\|q_i-p_i\|},\qquad
 0\leq h_C(\nu_i)-\langle \nu_i,p_i-z_c\rangle\leq\sigma,
 \qquad \sigma+\alpha\leq\varepsilon/16.
\]
Retain $\alpha\leq\gamma/2$ for the unsafe anchor. In the proof, the lower bound uses only $p_i\in C$; the upper bound becomes
$\langle u_i,z-p_i\rangle\leq3R(\alpha+\sigma)$.
Thus $\dist(z,C)\leq\ell^2/(2R)+6(\alpha+\sigma)\leq\varepsilon/2$ and the same component and margin guarantees hold. The gap is certifiable using a support upper bound and the feasible point, so no uncontrolled numerical normal error is assumed away. Exact projections give $\sigma=0$. For compact semialgebraic data, the stated real-algebraic optimization framework supplies exact projections; alternatively an oracle may return the certified feasible points and gap bounds above.

\section{Recovering the matrix from finite information}\label{sec:lp}
Assume also that $Z\subseteq z_c+R_Z\ball$ is certified. Put
\[
 M_B=R_Z+R+\alpha,\qquad \Lambda=M_B/\varepsilon.
\]
For the facets just constructed, the row multipliers with nominal margin $\varepsilon$ can be chosen to have row sums at most $\Lambda$. In fact, at any initial graph pair $(s_*,t_*)$ we have $\beta_j(s_*)\leq-\varepsilon$ and $\beta_i(t_*)\geq-M_B$. Hence every such row satisfies
\[
 -M_B+\varepsilon\sum_j A_{ij}
 \leq\beta_i(t_*)-\sum_j A_{ij}\beta_j(s_*)\leq-\varepsilon,
 \qquad \sum_jA_{ij}\leq(M_B-\varepsilon)/\varepsilon<\Lambda.
\]
A smaller row cap is permitted if a separate bound proves that nominal dual rows with margin $\varepsilon$ satisfy it.

Let $(s_\ell,t_\ell)=(\phi_d(x_\ell),\phi_d(f(x_\ell)))$ be \emph{exact graph pairs} at finitely many $x_\ell\in X$. Suppose a certified paired cover gives, for every $x\in X$, a common index $\ell$ with
\begin{equation}\label{eq:paircover}
 \|\phi_d(x)-s_\ell\|\leq\xi_s,\qquad
 \|\phi_d(f(x))-t_\ell\|\leq\xi_t.
\end{equation}
These are errors in the original features, before applying the facets. Choose a target propagation margin $\eta_*>0$ satisfying
\begin{equation}\label{eq:budget}
 \xi_t+\Lambda\xi_s+\eta_*\leq\varepsilon.
\end{equation}
For each row $i$, solve the following \emph{finite linear feasibility problem}:
\begin{equation}\label{eq:lp}
 \begin{aligned}
 A_{ij}&\geq0,\qquad \sum_jA_{ij}\leq\Lambda,\\
 \beta_i(t_\ell)-\sum_jA_{ij}\beta_j(s_\ell)
       +\xi_t+\xi_s\sum_jA_{ij}&\leq-\eta_*
       \qquad\text{for every }\ell.
 \end{aligned}
\end{equation}

\begin{theorem}[Feasible finite LP with continuum soundness]\label{th:lp}
Under the certified body, facet, row-cap, and exact paired-cover assumptions above, \eqref{eq:lp} is feasible. Every feasible solution gives
\[
 B\circ f\leq AB-\eta_*\one\quad\text{on all of }X,\qquad B_i=\beta_i\circ\phi_d,
\]
retaining the certified initial and unsafe margins. This inequality also holds on the full convexified relation $\F$.
\end{theorem}
\begin{proof}
Nominal dual rows with margin $\varepsilon$ satisfy every exact node constraint with residual at most $-\varepsilon$. Their row caps and \eqref{eq:budget} prove LP feasibility. For any returned row and any $x$, choose the common node in \eqref{eq:paircover}. Unit facet normals and $A_{ij}\geq0$ bound the change in the residual by
$\xi_t+\xi_s\sum_jA_{ij}$. The node constraint therefore implies the desired inequality at $x$. Affinity in the paired features extends it to convex combinations, and continuity to limits. Initial and unsafe separation were already certified geometrically.
\end{proof}

This LP does not impose a nonconvex spectral constraint: strict propagation and initial separation imply $\rho(A)<1$ by Theorem~\ref{th:main}. Exact rational LP data admit a rational feasible solution. Floating-point discovery requires a separate exact or interval verification of \eqref{eq:lp}; a solver's success status alone is insufficient. If node evaluations themselves are approximate, their certified errors must also be charged to the budget. The feasibility proof as written assumes exact graph nodes.

\textbf{Effective inputs and termination.}
A conditional construction takes as input a certificate of \eqref{eq:bodydata}, geometry bounds, certified support-value oracles, and a finite paired-cover oracle. A maximal net takes finitely many selections by Lemma~\ref{lem:net}; finitely many support calls and finitely many LPs then finish the construction. Proposition~\ref{prop:projection} substitutes finitely many exact projection calls for its sharper count. The procedure does not discover $C$ automatically.

One explicit admissible representation is compact semialgebraic $X,X_0,X_u$ and $C$, polynomial $f$, and algebraic coefficients, with certified algebraic feature coordinates and bounds. Real-algebraic decision/optimization oracles can verify the body conditions, separator, and support bounds. By Carath\'eodory's theorem, a pair in $\F\subseteq E\times E$ needs at most $2r+1$ graph atoms, so its use in the transition-inclusion formula is finite. Sphere-net selection is also a semialgebraic feasibility query. On a known bounding box, certified Lipschitz bounds $L_s$ for $\phi_d$ and $L_t$ for $\phi_d\circ f$ give \eqref{eq:paircover} from a state $h$-net with $\xi_s=L_sh$, $\xi_t=L_th$. A finite box subdivision, exact nonempty-cell decisions, and one exact point of $X$ per occupied cell yield such a cover. For guaranteed positive mesh size, choose $0<\eta_*<\varepsilon$ (for example $\eta_*=\varepsilon/2$). If $L_t+\Lambda L_s>0$, choosing
\[
 h\leq\frac{\varepsilon-\eta_*}{L_t+\Lambda L_s}
\]
suffices. The zero-denominator case needs only exact constant feature pairs. Algebraic LPs may be solved with exact algebraic arithmetic; rational data allow rational LPs. These statements give finite termination with the stated oracles, not efficiency or a general compact-continuous algorithm. In the worked example all body and cover conditions have direct analytic certificates and all LP checks are rational.

\section{Nonlinear certificate extraction and its scalar baseline}\label{sec:nonlinear}
Take
\[
 X=[-1,1],\quad f(x)=\tfrac12x+\tfrac14x^2,\quad
 X_0=[-\tfrac14,\tfrac14],\quad X_u=[\tfrac7{10},\tfrac45],\quad d=2.
\]
Since $f'(x)=(1+x)/2\in[0,1]$, $f(X)=[-1/4,3/4]\subseteq X$.
\textbf{Scalar baseline (pointed out by Majid Zamani).}
A single affine function already certifies this system:
\[
 b(x)=x-\tfrac35,\qquad A_{\rm sc}=[\tfrac12],\qquad
 b(f(x))-\tfrac12 b(x)=\tfrac14x^2-\tfrac3{10}\leq-\tfrac1{20}.
\]
Its initial and unsafe margins are respectively $7/20$ and $1/10$.
It belongs to the degree-at-most-two class and even has degree one.
Consequently the twelve-function construction below validates conditional extraction from a supplied certified body, but gives no advantage in function count, degree, or propagation margin over this scalar certificate. Its largest uniform propagation margin $73/2500$ for the fixed twelve-function pair is smaller than the scalar margin $1/20$. Section~\ref{sec:adding} supplies a separate nonlinear example with a proved benefit from adding functions.

\textbf{Supplied body and certified extraction data.}

Use $\phi_2(x)=(x,x^2)$ and $z_c=0$. For \emph{every} pair $(s,t)\in\mathcal F_2$,
\[
 t_1=\tfrac12s_1+\tfrac14s_2,\quad
 0\leq t_2\leq\tfrac9{16}s_2,\quad s_2\geq0.
\]
These follow by averaging the graph equalities and the pointwise bound
$f(x)^2=x^2(1/2+x/4)^2\leq(9/16)x^2$.
Consequently,
\[
 \|t\|_2^2\leq
 \left(\tfrac14+\tfrac1{16}+\tfrac{81}{256}\right)\|s\|_2^2
 =\tfrac{161}{256}\|s\|_2^2
 \leq(\tfrac{13}{16})^2\|s\|_2^2.
\]
This is a bound on mixed transition pairs, not just individual trajectories.

Let $C=(3/5)\mathbb B_2$, $\varepsilon=1/20$, and $\gamma=1/10$, with anchor $a_1=(1,0)$. On $I_2$,
$\|s\|_2\leq5/16$, so $I_2+\varepsilon\mathbb B_2\subseteq C$.
For any perturbed source $y=s+e\in C$ and independent successor error $w$, with $\|e\|_2,\|w\|_2\leq\varepsilon$,
\[
 \|t+w\|_2\leq\tfrac{13}{16}(3/5+\varepsilon)+\varepsilon
 =\tfrac{37}{64}<\tfrac35,
\]
leaving the certified inclusion gap $7/320$. The anchor strictly separates $J_2$ by $\gamma$. Thus the full robust convexified hypothesis holds. Notice that $f(1)=3/4\in X_u$; global avoidance in a single step is not assumed.

Here $r=2$ and $R/\varepsilon=12$. The direct support-net bound is $m\leq1537$, while Proposition~\ref{prop:projection} improves it to $m\leq112$: the expression inside its ceiling is $64\sqrt3$, between $110$ and $111$. Projection onto this ball is explicit. A smaller certified construction uses twelve rational unit normals, all with exact offset $3/5$:
\[
 \begin{gathered}
 a_1=(1,0),\ a_2=(0,1),\ a_3=(-1,0),\ a_4=(0,-1),\\
 a_5=(3/5,4/5),\ a_6=(3/5,-4/5),\
 a_7=(-3/5,4/5),\ a_8=(-3/5,-4/5),\\
 a_9=(4/5,3/5),\ a_{10}=(4/5,-3/5),\
 a_{11}=(-4/5,3/5),\ a_{12}=(-4/5,-3/5).
 \end{gathered}
\]
For $P=\{a_i^\top z\leq3/5\}$, exact two-facet vertex enumeration gives
$\max_{z\in P}\|z\|_2^2=2/5<(13/20)^2$, so \eqref{eq:otherpoly} holds. This is a direct geometric certificate; these twelve normals are not claimed to be the much finer universal net.

At the graph pair $(0,0)$, all facets equal $-3/5$, so nominal dual rows with margin $\varepsilon$ have row sums at most $11/12$. We may therefore take the sharper LP cap $\Lambda=1$.
Use the 161 exact nodes $x_\ell=-1+\ell/80$, $\ell=0,\ldots,160$.
Their covering radius is $h=1/160$. The derivative estimates
\[
 \|\phi_2'(x)\|_2^2=1+4x^2\leq5<3^2,\qquad
 \|(\phi_2\circ f)'(x)\|_2^2
 =f'(x)^2(1+4f(x)^2)\leq13/4<2^2
\]
give $\xi_s=3/160$, $\xi_t=1/80$, and
$\xi_t+\Lambda\xi_s=1/32$.
LP discovery used target $1/80$; exact replay accepts the target
$\eta_*=1/100$ with minimum LP slack $1/400$.

Specifically, for $B_i(x)=a_{i1}x+a_{i2}x^2-3/5$, take the matrix with only
\[
 A_{1,5}=28/125,\qquad A_{5,5}=52/125,\qquad A_{9,5}=58/125
\]
nonzero. The finite constraints \eqref{eq:lp} hold exactly at all 161 nodes. Theorem~\ref{th:lp} certifies the global continuum inequalities, including every convexified transition pair, with
\[
 \delta_0\geq\tfrac{23}{80},\qquad
 \delta_u\geq\tfrac1{10},\qquad
 \eta\geq\tfrac1{100},\qquad
 \rho(A)\leq\|A\|_\infty=\tfrac{58}{125}<1.
\]
This example verifies finite certificate extraction from the supplied certified body and cover. The scalar baseline above establishes that multiple components and degree two are unnecessary for safety of this particular system.

\textbf{Independent exact continuum check.}
For the same fixed $B,A$, expand the twelve polynomials
$r_i(x)=(AB(x))_i-B_i(f(x))$ in the degree-four Bernstein basis on $[-1,1]$:
\[
 r_i(2u-1)=\sum_{k=0}^4 c_{ik}\binom4k u^k(1-u)^{4-k},
 \qquad 0\leq u\leq1.
\]
Exact rational conversion gives $\min_{i,k}c_{ik}=73/2500$. The basis functions are nonnegative and sum to one, so $r_i(x)\geq73/2500$ on the entire interval, independently of the cover/Lipschitz argument. Moreover $r_1(1)=73/2500$; this is the exact largest uniform global propagation margin for this \emph{fixed} certificate pair. It is not an optimum over other functions or matrices. The LP guarantee remains the conservative $1/100$ after charging the feature-cover errors. The replay records all sixty Bernstein coefficients and checks both proofs.

\section{Nonlinear recovery by adding functions at fixed degree}\label{sec:adding}
We strengthen the original contracting-rotation example to a uniform neighborhood of nonlinear maps. Write
\[
 X=\{x\in\R^2:\|x\|_2\leq1\},\quad
 X_0=\{x:\|x\|_2\leq\tfrac14\},\quad
 X_u=\{(\tfrac34,0)\},\quad
 R=\begin{pmatrix}0&-1\\1&0\end{pmatrix},\quad q=\tfrac9{10}.
\]
\begin{proposition}[A perturbation-stable benefit from more affine functions]\label{prop:nonlinear}
Let $f:X\to\R^2$ be continuous and satisfy
\[
 f(x)=qRx+e(x),\qquad \sup_{x\in X}\|e(x)\|_2\leq\delta,\qquad 0\leq\delta<\tfrac1{20}.
\]
Then $f(X)\subseteq X$. At prescribed degree $d=1$, no scalar ordinary comparison certificate exists for any multiplier $\lambda\geq0$, whereas four affine components give a strict fixed-anchor VBC. The latter conclusion holds on the full convexified graph.
\end{proposition}
\begin{proof}
First, $\|f(x)\|_2\leq q+\delta<1$ on $X$.
Suppose a scalar affine certificate is $b(x)=a^\top x-c$.
Strict negativity at $0\in X_0$ gives $c>0$, and the unsafe condition implies
\[
 c<\tfrac34a_1\leq\tfrac34\|a\|_2,\qquad a\neq0.
\]
For $x_*=R^\top a/\|a\|_2\in X$, skew symmetry gives $a^\top x_*=0$, and hence
\[
 b(x_*)=-c<0,\qquad
 b(f(x_*))\geq(q-\delta)\|a\|_2-c>0,
\]
because $q-\delta>17/20>3/4$. This contradicts
$b\circ f\leq\lambda b$ for every $\lambda\geq0$.

For the four-component certificate put $r_0=1/2$ and
\[
 G=\begin{pmatrix}1&0\\0&1\\-1&0\\0&-1\end{pmatrix},\quad
 \Pi=\begin{pmatrix}0&0&0&1\\1&0&0&0\\0&1&0&0\\0&0&1&0\end{pmatrix},
 \qquad B(x)=Gx-r_0\one,\quad A=q\Pi.
\]
Since $GR=\Pi G$, direct substitution yields
\[
 B(f(x))-AB(x)=Ge(x)-(1-q)r_0\one
 \leq-\bigl(\tfrac1{20}-\delta\bigr)\one.
\]
Every component is at most $-1/4$ on $X_0$, $B_1(3/4,0)=1/4$,
and $\rho(A)=q<1$. Affinity in $(s,t)$ extends the comparison from the graph to its full convex hull.
\end{proof}

\textbf{An explicit cubic family.}
For every $0<\theta\leq1/40$ take
\begin{equation}\label{eq:cubic-family}
 f_\theta(x)=\left[\tfrac9{10}-\theta(1-\|x\|_2^2)\right]Rx.
\end{equation}
Its degree is three, while the prescribed certificate degree remains one.
Here $e(x)=-\theta(1-\|x\|_2^2)Rx$ has norm at most $\theta$ on $X$,
so Proposition~\ref{prop:nonlinear} applies uniformly with $\delta=1/40$.
The same four affine functions and matrix certify every member of this family with
\[
 \delta_0=\delta_u=\tfrac14,\qquad \eta=\tfrac1{40},\qquad \rho(A)=\tfrac9{10}.
\]
This proves an actual nonlinear same-degree recovery. Four is an upper bound, not a proved minimum component count. The uniform perturbation proposition also shows that the recovery does not rely on one isolated polynomial identity.

\textbf{Increasing degree is an alternative.}
For the family \eqref{eq:cubic-family}, the scalar quadratic
\[
 b_2(x)=\|x\|_2^2-\tfrac14,\qquad A_2=[\tfrac{81}{100}]
\]
satisfies
\[
 b_2(f_\theta(x))-\tfrac{81}{100}b_2(x)
 =\left(\left[\tfrac9{10}-\theta(1-\|x\|_2^2)\right]^2-\tfrac{81}{100}\right)
       \|x\|_2^2-\tfrac{19}{400}
 \leq-\tfrac{19}{400}.
\]
Its initial and unsafe margins are $3/16$ and $5/16$. Thus the minimum
scalar forward constant-comparison degree is exactly two, while the minimum degree
allowing a finite fixed-anchor VBC is exactly one: degree zero cannot separate
nonempty initial and unsafe sets. At degree one one must add functions;
increasing to degree two permits one function.

\textbf{The chosen-body hypothesis, including mixtures and both errors.}
For the cubic family use the uniform bound $\delta=1/40$. At $d=1$, a pair $(s,t)\in\mathcal F_1$ has
\[
 t=qRs+\bar e,\qquad\|\bar e\|_2\leq\delta,
\]
by averaging $e(x)$; it need not satisfy $t=f(s)$.
Set $C=[-1/2,1/2]^2$ and $\varepsilon=1/100$.
For every such pair, if $y=s+v\in C$ and
$\|v\|_2,\|w\|_2\leq\varepsilon$, then
\[
 \|t+w\|_\infty
 \leq qr_0+\delta+(q+1)\varepsilon
 \leq\tfrac{247}{500}<\tfrac12
 \qquad(\delta=1/40).
\]
Also $I_1+\varepsilon\ball\subseteq C$ and $C\cap J_1=\varnothing$.
This proves $T_\varepsilon(C)\subseteq C$ for all convexified transition pairs
and independent source/successor errors. Subsequent convex combinations remain in $C$.
The original linear rotation is the limiting case $e=0$, with the sharper
propagation margin $1/20$ and robust square bound $469/1000$.

\section{An obstruction: adding same-degree functions cannot succeed}
Reuse the nonlinear example from the earlier note:
\begin{align*}
 f(x)&=\tfrac78x+\tfrac58x^3-\tfrac12x^5,\qquad X=[-1,1],\\
 a&=\tfrac{12}{25},\quad u=\tfrac{1641}{3200},\quad v=\tfrac{103}{200},\\
 X_0&=[-a,a],\qquad X_u=[-v,-u]\cup[u,v].
\end{align*}
Writing $z=x^2$ gives
\[
 f'(x)=\tfrac14+(1-z)(\tfrac58+\tfrac52z)\geq\tfrac14\quad(x\in X);
\]
together with $f(\pm1)=\pm1$, this proves that $f$ is \emph{onto} $X$.
There is no true unsafe trajectory from $X_0$: the single quadratic
\[
 b(x)=x^2-\tfrac14
\]
is an \emph{interpolation-inspired barrier certificate (IBC)} with one frame. If $b(x)\leq0$, then $f(x)=xh(x^2)$, where
$h(z)=1+\tfrac12(z-\tfrac14)(1-z)$ lies in $[7/8,1]$ for $0\leq z\leq1/4$. Hence $|f(x)|\leq|x|\leq1/2$. The initial and unsafe margins are positive because $a<1/2<u$.

Nevertheless, at degree two put
\[
 \mu=\tfrac{369}{625}\delta_0+
     \tfrac{128}{625}(\delta_{-3/4}+\delta_{3/4}),\qquad
 \mu_0=\tfrac12(\delta_{-a}+\delta_a),\qquad
 \nu=\tfrac12(\delta_{-u}+\delta_u).
\]
All are probability measures with support in their designated sets. Since
$f(3/4)=1641/2048$ and $u=a f(3/4)/(3/4)$,
\begin{equation}\label{eq:witness}
 \int\phi_2\dd\mu=\int\phi_2\dd\mu_0=(0,a^2),\qquad
 \int\phi_2\circ f\dd\mu=\int\phi_2\dd\nu=(0,u^2).
\end{equation}
Thus $\bigl((0,a^2),(0,u^2)\bigr)\in\mathcal F_2$ has source in $I_2$ and successor in $J_2$. Already $C_0^1$ meets $J_2$, and so does $\mathcal K_\varepsilon$ for every $\varepsilon\geq0$.

\begin{proposition}[No finite same-degree repair]\label{prop:obstruction}
For this example no forward fixed-anchor VBC of degree at most two exists, with any finite component count and any constant $A\geq0$. The conclusion also holds through degree three.
\end{proposition}
\begin{proof}
For any degree-at-most-two vector $B$, \eqref{eq:witness} and the unit masses imply
\[
 \int B\circ f\dd\mu=\int B\dd\nu
 \leq A\int B\dd\mu
 =A\int B\dd\mu_0\leq0.
\]
The designated unsafe component has integral at least $\delta_u>0$, a contradiction. All odd moments through degree three vanish by symmetry, so the same identities and argument apply to cubic components.
\end{proof}
The source atoms at $\pm3/4$ lie outside $X_0$. Their mixtures imitate initial features while the successor mixtures imitate unsafe features. This is precisely the information that ordinary trajectory safety omits. The fixed-anchor restriction is essential: positivity of some component at each unsafe point does not in general imply positivity of one component averaged over $\nu$.

\textbf{The quartic recovery belongs to the neutral boundary class.}
The earlier note gives the ordinary scalar quartic
\[
 p(x)=\tfrac3{10000}-(x^2-\tfrac14)^2,\qquad A=[1],
\]
with strict initial/unsafe separation and $p\circ f\leq p$ on $X$. Its exact verification is retained in the repository. Since this $f$ is onto, Corollary~\ref{cor:onto} excludes \emph{every} strict or Schur-stable VBC at every degree and component count. Thus the robust converse does not explain the quartic recovery. Every ordinary VBC here must have $\rho(A)=1$: the initial-set Perron argument gives $\rho(A)\leq1$, and onto-ness excludes $\rho(A)<1$.

\section{Physical robust safety does not imply feature robustness}\label{sec:physical}
Consider the nonidentity map
\[
 X=[-1,1],\quad f(x)=x^2,\quad
 X_0=[-1/4,1/4],\quad X_u=[3/4,1].
\]
The ordinary certificate $B(x)=x^2-1/2$, $A=[1]$ has
\[
 B\leq-7/16\ \text{on }X_0,\qquad B\geq1/16\ \text{on }X_u,\qquad
 B(f(x))-B(x)=x^4-x^2\leq0\ \text{on }X.
\]
Yet $1\in X_u$ is a fixed point. No strict or Schur-stable fixed-anchor VBC exists at any degree or finite component count. By Theorem~\ref{th:main}, the robust convexified-feature hypothesis fails at every prescribed degree.

In contrast, additive \emph{state} disturbances $|d_k|\leq1/8$ preserve the safe interval $C_{\rm state}=[-1/2,1/2]$, because
\[
 f(C_{\rm state})+[-1/8,1/8]=[-1/8,3/8]\subseteq C_{\rm state}.
\]
This interval contains $X_0$ and avoids $X_u$, proving physical robust safety. Independent source and successor feature errors can move off the realizable monomial variety; they encode a different and stronger requirement in this example.

The original identity example gives the same ordinary-versus-strict distinction at degree one: $f(x)=x$ on $X=[-1,1]$, $X_0=[-1/2,0]$, $X_u=[1/2,1]$, and $B(x)=x-1/4$, $A=[1]$ satisfy ordinary separation and propagation. Every $\varepsilon>0$ allows successive feature increments of size at most $\varepsilon$ from $0$ to $1/2$, so $\mathcal K_\varepsilon$ reaches $J_1$.

\section{Theorem-specific comparison with existing results}\label{sec:literature}
The comparison below separates inherited mechanisms from the precise claims
to be assessed for novelty. A difference in assumptions is not by itself a priority claim.

\subsection{Theorem~\ref{th:main}: the existence equivalence}
\textbf{Closest multiplier result: Horv\'ath--Song--Terlaky, Theorem 3.10 and Remark 3.11.}
For a supplied convex function $g$ and its sublevel set $\{g\leq0\}$, \cite{hst} assumes
strict feasibility and concavity of $g\circ f_d$, and characterizes its
invariance by $g\circ f_d\leq\lambda g$ globally on $\R^n$, for
$\lambda\geq0$. If both functions are quadratic, the S-lemma removes the
successor-concavity assumption; Remark 3.11 extends the first statement
to several convex inequalities.
Our Lemma~\ref{lem:dual} is the same classical duality mechanism on a compact
convex relation. Theorem~\ref{th:main} instead starts with no supplied
polynomial family, assumes positive-radius separation after full graph and
reachable-set convexification in a fixed feature space, and concludes existence
of finitely many degree-bounded functions with a designated unsafe anchor and
$\rho(A)<1$, together with the converse implication.
No concavity of $B_i\circ f$ is required, but ordinary state-space invariance
has been replaced by a stronger convexified-feature robustness condition.
Thus the theorem is not an unconditional removal of their concavity assumptions.

\textbf{Closest smooth safety converse: Liu, Section IV, Theorem 30.}
The discrete-time result in \cite{liu} uses additive state-disturbance robust
safety, separation of the disturbed reachable closure from the unsafe set,
local Lipschitz regularity, and compactness of that closure or global
Lipschitz continuity. It gives a smooth robust barrier at a smaller disturbance
radius. Our hypothesis uses independent feature errors and convexified
transition pairs; our conclusion prescribes the polynomial degree and a
constant nonnegative Schur comparison. Section~\ref{sec:physical} proves
that physical robustness alone need not satisfy this feature premise.
Neither converse subsumes the other through a simple change in disturbance radius.

\textbf{Closest polynomial-cone/positive-system result: Berger--Ghanbarpour--Sankaranarayanan, HSCC 2024.}
Theorem 3.5 of \cite{cone2024} proves invariance for polynomial ODEs from
$L_f(g)+\lambda g\in\operatorname{cone}(K G^{\leq\ell})$, using positive
systems. Definition 5.1 and Corollary 5.2 turn this sufficient condition into
a refinement fixed-point criterion. It permits polynomial multipliers and
products of generators, and already uses polynomial templates and multiple
inequalities. Our constant-matrix coupling is more restricted. Our equivalence concerns
robust convexified fixed-feature separation with an unsafe anchor and Schur scope.
Neither polynomial features, positive comparison systems, nor cone refinement
is claimed new here.

\subsection{Theorem~\ref{th:count} and Proposition~\ref{prop:projection}: conditional facets}
\textbf{Closest geometry: Har-Peled--Jones, Theorem 1.4.}
The classical Dudley result in \cite{harpeled} approximates a supplied bounded
convex body by $O_r((R/\tau)^{(r-1)/2})$ supporting halfspaces within
Hausdorff error $\tau$. Proposition~\ref{prop:projection} applies this mechanism
with a prescribed unsafe separator, explicit support/projection-gap error
budgets, and $\tau=\varepsilon/2$. The exponent is classical.
Theorem~\ref{th:count} gives a coarser direct support-net alternative.
The degree bound comes from composing affine facets with the fixed feature map,
not from a new convex approximation theorem.

\textbf{Closest dynamical approximation: Alessio--Lazar--Bemporad--Heemels, Theorem 4.}
Theorem 4 in \cite{alessio} is a general intermediate-set inclusion result for
perturbed nonlinear discrete-time dynamics: an approximation within a
contraction margin preserves the corresponding invariance/contraction property.
Ellipsoids enter its later finite construction, not the general theorem.
Our input is a supplied certified feature body satisfying \eqref{eq:bodydata};
the output additionally includes initial and fixed-anchor margins, a
prescribed-degree family, and a global comparison matrix.
The approximation-within-margin principle and multiple-successor uncertainty
are inherited mechanisms, not independent contributions.

\subsection{Theorem~\ref{th:lp}: conditional matrix extraction}
\textbf{Closest finite LP synthesis: Berger--Sankaranarayanan, Sections 6--7 and 9.}
Proposition 13 in \cite{berger} gives a convex candidate problem (an LP for
polyhedral norms), and Proposition 15 gives verification LPs; Proposition 19 and Theorem 21 prove
coverage-based progress and a packing termination bound.
Section 9 explicitly discusses discrete-time piecewise-linear extensions,
and future safety/invariant-polytope extensions.
Thus changing continuous time to discrete time, using finite LPs, or deriving
a packing bound is insufficient as a novelty claim.
Theorem~\ref{th:lp} assumes a supplied certified body and selected facets,
a row cap and exact paired feature cover. It guarantees feasibility and
continuum validity of a fixed-anchor comparison matrix after charging both
source and successor errors. It does not discover the body or adaptively
discover a Lyapunov function.

\textbf{Polynomial refinement is also prior art.}
Algorithms 2--6 and Lemmas 5.4, 5.5, and 5.8 of \cite{cone2024} already address
finite polynomial-cone refinement, robust projections, and widening for
termination. Widening may return a trivial invariant.
Our conditional extraction starts after a useful robust body is supplied and
guarantees a strictly separating certificate under its stated data budgets.
The projection in Dudley's feature-space geometry and their projection in
polynomial-coefficient cone space solve different subproblems.
This comparison establishes no direct implication to our fixed-feature/Schur
equivalence; its publication-level priority remains to be assessed.

\section{Assessment, remaining questions, and reproducibility}
The Schur characterization, dynamical limitations, and physical-robustness counterexample formalize observations supplied in Majid Zamani's comments. The quantitative bounds use classical support-net and Dudley geometry; the row LP extracts a certificate from supplied, certified body and cover data. The October 23 response adds the scalar baseline, a perturbation-stable nonlinear recovery proposition, an exact scalar/vector degree tradeoff, and a consolidated manuscript outline.

The revision proves the spectral characterization requested in the comments, makes its stronger dynamical limitations explicit, and adds constructive consequences: the bounds \eqref{eq:mcount} and \eqref{eq:projection-count}, and the finite feasible LP \eqref{eq:lp}. The scalar nonlinear example validates conditional extraction from its supplied body; its affine baseline is simpler. Proposition~\ref{prop:nonlinear} separately proves a nonlinear same-degree function-count benefit, and \eqref{eq:cubic-family} gives an exact scalar/vector degree tradeoff. The general theorems are mathematical proofs in this note; their formal machine verification is not claimed.

The ordinary boundary class $\rho(A)=1$ still lacks a general converse here; onto examples prevent recovering it by imposing strictness. As requested, solving that boundary case is not a prerequisite for drafting the consolidated paper. Automatic discovery of the body $C$ and minimal component counts also remain open. The projection rate is classical, and publication-level novelty of the combined feature-space equivalence and certified extraction requires further comparison with the closest converse and invariant-approximation results.

\subsection*{Consolidated manuscript outline}
Working title: \emph{Structural Conservatism of Discrete-Time Barrier Certificates:
Moment Obstructions, a Schur-Stable Converse, and Exact Degree Separations}.
The paper can now be drafted around the following sequence.
\begin{enumerate}
\item \textbf{Problem and certificate classes.}
Fix degree, distinguish scalar/vector comparisons from implication-style
IBCs, and state the fixed-anchor and propagation-direction restrictions.
Separate ordinary $\rho(A)=1$ certificates from the Schur subclass.
\item \textbf{Convexified comparisons and moment obstructions.}
Present the fixed-family dual characterization, then the finite-support
moment criterion excluding \emph{every finite component count} at degree $d$.
Explain that a witness describes incompatible averages, not an unsafe trajectory.
State its sufficient, rather than necessary, scope.
\item \textbf{Schur-stable existence converse.}
Prove Theorem~\ref{th:main}, treating both convexifications and independent
feature errors. Derive uniform eventual avoidance and the onto/periodic
obstructions. Keep this forward-only characterization separate from the
earlier forward/backward moment obstruction.
\item \textbf{Conditional finite extraction.}
Starting from a supplied certified body $C$, give the attributed facet
approximation, explicit margins/counts, and paired-cover row LP.
Separate existence, conditional finite extraction, and discovery.
\item \textbf{When more functions help.}
Use Proposition~\ref{prop:nonlinear} and the cubic family: one affine function
fails, four succeed, and one quadratic succeeds. Thus function count and
degree provide distinct ways to remove this obstruction.
\item \textbf{When no number of same-degree functions helps.}
Use the exact nonlinear moment witness to exclude all degree-at-most-three
fixed-anchor comparison VBCs, while a quadratic single-frame
interpolation-inspired certificate and a quartic ordinary comparison work.
The minima two and four concern \emph{different propagation classes}.
The quartic matrix has $\rho(A)=1$ and lies outside the Schur converse.
\item \textbf{Evidence, comparisons, and implications.}
Include the twelve-function extraction example with its scalar baseline,
exact rational replays, and theorem-specific prior-work comparisons.
Conclude only the proved formulation-selection implications: try more functions
where a construction permits it; a moment obstruction rules out that remedy
and requires changing degree or certificate conditions within the stated classes.
Do not infer a complete decision procedure from these sufficient tests.
\end{enumerate}
The four proof pillars are therefore the component-count-independent
obstruction, the Schur-stable converse, conditional extraction, and exact
degree separations. The neutral-class converse is a stated open direction,
not a missing prerequisite for this manuscript plan.

\textbf{Code and reproducibility.}
All computational materials are in the public Certificate Expressiveness repository:
\begin{center}
\url{https://github.com/rezairaji60/certificate-expressiveness}
\end{center}
The Overleaf-ready source is \texttt{paper/converse.tex}; the exact nonlinear matrix and its deterministic report are
\texttt{evidence/nonlinear\_comparison\_matrix.json} and
\texttt{evidence/constructive\_converse\_report.json}.
Run \texttt{python src/constructive\_converse.py} for a standard-library-only exact replay.
Optional \texttt{-{}-discover} uses SciPy solely to propose a matrix, which is accepted only after rational verification.
The existing \texttt{src/exact\_verify.py} and
\texttt{src/exact\_verify\_converse.py} retain the original continuum, rotation, and moment checks. The latter also replays the scalar baseline, cubic-family identities, nonlinear four-function margins, and mixed-transition robustness bound. Tests reject corrupted comparison matrices and understated feature-cover errors as well as invalid moment witnesses.
The finite cover's soundness rests on the stated analytic derivative bounds; finite node checks without those bounds would not certify the continuum. The nonlinear example additionally has an independent exact Bernstein proof and an endpoint witness for its fixed-certificate propagation margin.

\clearpage
\begin{thebibliography}{6}
\small\raggedright
\bibitem{hst}
Z. Horv\'ath, Y. Song, and T. Terlaky,
\emph{Invariance Conditions for Nonlinear Dynamical Systems},
arXiv:1607.01107v1, 2016. Theorem 3.10 and Remark 3.11.
\par\url{https://arxiv.org/abs/1607.01107}.
\bibitem{liu}
J. Liu, \emph{Converse Barrier Functions via Lyapunov Functions},
IEEE Transactions on Automatic Control, vol. 67, no. 1, pp. 497--503, 2022.
DOI: 10.1109/TAC.2021.3085419.
\par\url{https://arxiv.org/abs/2007.11086}.
\bibitem{alessio}
A. Alessio, M. Lazar, A. Bemporad, and W. P. M. H. Heemels,
\emph{Squaring the circle: An algorithm for generating polyhedral invariant sets from ellipsoidal ones},
Automatica, vol. 43, no. 12, pp. 2096--2103, 2007.
DOI: 10.1016/j.automatica.2007.04.028.
\par\url{https://cse.lab.imtlucca.it/~bemporad/publications/papers/automatica_squaring_circle.pdf}.
\bibitem{berger}
G. O. Berger and S. Sankaranarayanan,
\emph{Counterexample-guided computation of polyhedral Lyapunov functions for piecewise linear systems},
Automatica, vol. 155, article 111165, 2023.
DOI: 10.1016/j.automatica.2023.111165.
\bibitem{cone2024}
G. O. Berger, M. Ghanbarpour, and S. Sankaranarayanan,
\emph{Cone-Based Abstract Interpretation for Nonlinear Positive Invariant Synthesis},
Proceedings of HSCC '24, pp. 1--16, 2024.
DOI: 10.1145/3641513.3650127.
\par\url{https://guberger.github.io/assets/papers/10.1145_3641513.3650127.pdf}.
\bibitem{harpeled}
S. Har-Peled and M. Jones,
\emph{A Proof of Dudley's Convex Approximation},
arXiv:1912.01977v3, 2024. Theorem 1.4.
\par\url{https://arxiv.org/abs/1912.01977v3}.
\end{thebibliography}
\end{document}
